% =============================================================================
% 3.2 Infrastructure Strategy — Local-First with Optional Cloud Validation
% =============================================================================

\section{Infrastructure Strategy --- Local-First with Optional Cloud Validation}
\label{sec:infrastructure}

A critical design decision in this thesis is the \textbf{local-first infrastructure strategy}, which addresses budget uncertainty and maximises experimental control.

\subsection{Primary Environment: Local Infrastructure}
\label{subsec:local-infra}

All experiments are designed to run entirely on the researcher's local machine using:

\begin{description}[style=nextline]
  \item[VM Baseline] A Linux virtual machine provisioned via \textbf{Multipass} (Ubuntu 22.04), configured with 4~vCPU and 4~GB RAM. Dagster is deployed with the \texttt{DockerRunLauncher}, which launches each pipeline run as an isolated Docker container on the VM.
  \item[Kubernetes Comparison] A \textbf{Kind} (Kubernetes in Docker) cluster running on the same physical host, configured with resource limits matching the VM. Dagster is deployed with the Kubernetes Run Launcher via Helm.
\end{description}

This approach provides:

\begin{itemize}
  \item \textbf{Zero cloud cost}---no budget dependency or risk of mid-experiment credit exhaustion.
  \item \textbf{Full control}---exact CPU/memory allocation is enforced and reproducible.
  \item \textbf{Faster iteration}---no cloud provisioning delays between experiment runs.
  \item \textbf{Reproducibility}---any researcher with a comparable machine can replicate the setup.
\end{itemize}

The shared physical host is an intentional design choice: it ensures both environments have access to the same underlying hardware, making the comparison about \emph{execution model architecture} (VM Docker containers vs.\ Kubernetes pods) rather than about hardware differences.

\subsection{Optional Cloud Validation Phase}
\label{subsec:cloud-validation}

If GCP budget permits (free tier credits or education grants), a subset of experiments will be repeated on:

\begin{description}
  \item[VM] Google Compute Engine instance (e2-standard-2: 2~vCPU, 8~GB RAM).
  \item[Kubernetes] Google Kubernetes Engine (GKE) standard cluster with Kubernetes Run Launcher.
\end{description}

Cloud results serve as a \textbf{directional validation} of local findings, not as the primary dataset. If cloud validation is not possible, this is documented as a delimitation, and the thesis argues that the architectural comparison (shared vs.\ isolated execution) is valid regardless of whether the underlying host is a laptop or a cloud instance---a position supported by the infrastructure-agnostic nature of resource contention theory \citep{nanda1991}.

\subsection{Infrastructure Specifications}
\label{subsec:infra-specs}

\begin{table}[H]
  \centering
  \caption{Infrastructure comparison between VM and Kubernetes environments.}
  \label{tab:infra-specs}
  \small
  \begin{tabularx}{\textwidth}{lXX}
    \toprule
    \textbf{Component} & \textbf{VM Baseline} & \textbf{Kubernetes (Kind)} \\
    \midrule
    Provisioning        & Multipass (Ubuntu 22.04)       & Kind v0.20+ on Docker \\
    CPU                 & 4 vCPU (fixed)                 & 4 vCPU (node resource limit) \\
    Memory              & 4 GB RAM (fixed)               & 8 GB RAM (node resource limit) \\
    Dagster executor    & DockerRunLauncher              & K8sRunLauncher via Helm \\
    Resource isolation  & Per-container (shared host kernel OOM) & Per-pod CPU/memory requests and limits \\
    Autoscaling         & Not applicable                 & No autoscaling (single-node Kind cluster) \\
    IaC                 & Shell script / Multipass config & Kind config YAML + Helm \\
    \bottomrule
  \end{tabularx}
\end{table}
